% GENERATED FILE. Edit canonical lessons, then run npm run generate:books.
% canonical-source-hash: fnv1a64:07ff66c8fc8a7e22
% canonical-lessons: JA-C150-uwagi, JA-C150-tebukuro, JA-C150-hazusu, JA-C150-musubu, JA-C150-yubiwa

\chapter{Jackets, Gloves and Rings}
\label{ch:ja-jackets-gloves-and-rings}
\hlchaptermodality{150}

\begin{chapteropening}
\textbf{By the end of this chapter:} \emph{I can talk about jackets, gloves and rings.}

Complete the last lesson of chapter 150: I can talk about jackets, gloves and rings.
\end{chapteropening}

\section[\texorpdfstring{\ja{うわぎ} (uwagi)}{うわぎ (uwagi)}]{\ja{うわぎ} (uwagi) — a jacket; a coat}

\label{lesson:JA-C150-uwagi}



\begin{quote}
Before the new one: say the Japanese for Western clothes, then the Japanese for underwear.
\end{quote}

\subsection*{You'll want to know: \ja{うわぎ}}
\textbf{\ja{うわぎ}} — \emph{uwagi} — \textquotedblleft{}a jacket; a coat\textquotedblright{}.

\begin{quote}
\textbf{\ja{ふゆは} \ja{うわぎを} \ja{きます。}} — \emph{fuyu wa uwagi o kimasu} — \textquotedblleft{}In winter I wear a jacket.\textquotedblright{}
\end{quote}

\begin{cousinweb}
\textbf{\ja{うわぎ}} (\emph{uwagi}) is native Japanese, built from \emph{uwa-}, \textquotedblleft{}upper\textquotedblright{}, the joining form of \emph{ue}, and \emph{ki}, \textquotedblleft{}wearing\textquotedblright{}, voiced to \emph{gi}; both parts are inherited rather than borrowed (JMdict, KANJIDIC2).
\end{cousinweb}

\begin{grammarlens}[title={What you've built}]
One more word for everyday life at A2.
\end{grammarlens}

\subsection*{Your turn}
\begin{itemize}
\raggedright
  \item \emph{Say it:} \emph{uwagi}
  \item \emph{Say it:} \emph{uwagi}, once more
  \item \emph{Say it:} \emph{fuyu wa uwagi o kimasu} — In winter I wear a jacket.
  \item \emph{Recall:} say the Japanese for Western clothes, then the Japanese for underwear, then say \emph{uwagi} again
\end{itemize}

\begin{tcolorbox}[breakable,colback=teal!4,colframe=teal!35!black,title={Before you move on}]
What does \ja{うわぎ} mean? (\textquotedblleft{}A jacket; a coat\textquotedblright{}.) Say it once more.
\end{tcolorbox}
\section[\texorpdfstring{\ja{てぶくろ} (tebukuro)}{てぶくろ (tebukuro)}]{\ja{てぶくろ} (tebukuro) — gloves}

\label{lesson:JA-C150-tebukuro}



\begin{quote}
Before the new one: say the Japanese for underwear, then the Japanese for a jacket.
\end{quote}

\subsection*{You'll want to know: \ja{てぶくろ}}
\textbf{\ja{てぶくろ}} — \emph{tebukuro} — \textquotedblleft{}gloves\textquotedblright{}.

\begin{quote}
\textbf{\ja{ふゆは} \ja{てぶくろを} \ja{つかいます。}} — \emph{fuyu wa tebukuro o tsukaimasu} — \textquotedblleft{}In winter I use gloves.\textquotedblright{}
\end{quote}

\begin{cousinweb}
\textbf{\ja{てぶくろ}} (\emph{tebukuro}) is native Japanese, built from \emph{te}, \textquotedblleft{}hand\textquotedblright{}, and \emph{fukuro}, \textquotedblleft{}bag\textquotedblright{}, voiced to \emph{bukuro}; both parts are inherited rather than borrowed (JMdict, KANJIDIC2).
\end{cousinweb}

\begin{grammarlens}[title={What you've built}]
One more word for everyday life at A2.
\end{grammarlens}

\subsection*{Your turn}
\begin{itemize}
\raggedright
  \item \emph{Say it:} \emph{tebukuro}
  \item \emph{Say it:} \emph{tebukuro}, once more
  \item \emph{Say it:} \emph{fuyu wa tebukuro o tsukaimasu} — In winter I use gloves.
  \item \emph{Recall:} say the Japanese for underwear, then the Japanese for a jacket, then say \emph{tebukuro} again
\end{itemize}

\begin{tcolorbox}[breakable,colback=teal!4,colframe=teal!35!black,title={Before you move on}]
What does \ja{てぶくろ} mean? (\textquotedblleft{}Gloves\textquotedblright{}.) Say it once more.
\end{tcolorbox}
\section[\texorpdfstring{\ja{はずす} (hazusu)}{はずす (hazusu)}]{\ja{はずす} (hazusu) — to remove; to take off (glasses, a ring)}

\label{lesson:JA-C150-hazusu}



\begin{quote}
Before the new one: say the Japanese for a jacket, then the Japanese for gloves.
\end{quote}

\subsection*{You'll want to know: \ja{はずす}}
\textbf{\ja{はずす}} — \emph{hazusu} — \textquotedblleft{}to remove; to take off (glasses, a ring)\textquotedblright{}.

\begin{quote}
\textbf{\ja{うちで} \ja{めがねを} \ja{はずします。}} — \emph{uchi de megane o hazushimasu} — \textquotedblleft{}At home I take off my glasses.\textquotedblright{}
\end{quote}

Polite \textbf{\ja{はずします}}, \emph{hazushimasu}; a godan verb.

\begin{cousinweb}
\textbf{\ja{はずす}} (\emph{hazusu}) is a native Japanese word (wago), inherited rather than borrowed from Chinese; its kanji was attached later as a kun'yomi, and its older origin is not settled (JMdict, KANJIDIC2).
\end{cousinweb}

\begin{grammarlens}[title={What you've built}]
One more word for everyday life at A2.
\end{grammarlens}

\subsection*{Your turn}
\begin{itemize}
\raggedright
  \item \emph{Say it:} \emph{hazusu}
  \item \emph{Say it:} \emph{hazusu}, once more
  \item \emph{Say it:} \emph{uchi de megane o hazushimasu} — At home I take off my glasses.
  \item \emph{Recall:} say the Japanese for a jacket, then the Japanese for gloves, then say \emph{hazusu} again
\end{itemize}

\begin{tcolorbox}[breakable,colback=teal!4,colframe=teal!35!black,title={Before you move on}]
What does \ja{はずす} mean? (\textquotedblleft{}To remove; to take off (glasses, a ring)\textquotedblright{}.) Say it once more.
\end{tcolorbox}
\section[\texorpdfstring{\ja{むすぶ} (musubu)}{むすぶ (musubu)}]{\ja{むすぶ} (musubu) — to tie; to knot}

\label{lesson:JA-C150-musubu}



\begin{quote}
Before the new one: say the Japanese for gloves, then the Japanese for to remove.
\end{quote}

\subsection*{You'll want to know: \ja{むすぶ}}
\textbf{\ja{むすぶ}} — \emph{musubu} — \textquotedblleft{}to tie; to knot\textquotedblright{}.

\begin{quote}
\textbf{\ja{くつの} \ja{ひもを} \ja{むすびます。}} — \emph{kutsu no himo o musubimasu} — \textquotedblleft{}I tie my shoelaces.\textquotedblright{}
\end{quote}

Polite \textbf{\ja{むすびます}}, \emph{musubimasu}; a godan verb.

\begin{cousinweb}
\textbf{\ja{むすぶ}} (\emph{musubu}) is a native Japanese word (wago), inherited rather than borrowed from Chinese; its kanji was attached later as a kun'yomi, and its older origin is not settled (JMdict, KANJIDIC2).
\end{cousinweb}

\begin{grammarlens}[title={What you've built}]
One more word for everyday life at A2.
\end{grammarlens}

\subsection*{Your turn}
\begin{itemize}
\raggedright
  \item \emph{Say it:} \emph{musubu}
  \item \emph{Say it:} \emph{musubu}, once more
  \item \emph{Say it:} \emph{kutsu no himo o musubimasu} — I tie my shoelaces.
  \item \emph{Recall:} say the Japanese for gloves, then the Japanese for to remove, then say \emph{musubu} again
\end{itemize}

\begin{tcolorbox}[breakable,colback=teal!4,colframe=teal!35!black,title={Before you move on}]
What does \ja{むすぶ} mean? (\textquotedblleft{}To tie; to knot\textquotedblright{}.) Say it once more.
\end{tcolorbox}
\section[\texorpdfstring{\ja{ゆびわ} (yubiwa)}{ゆびわ (yubiwa)}]{\ja{ゆびわ} (yubiwa) — a ring (for a finger)}

\label{lesson:JA-C150-yubiwa}



\begin{quote}
Before the new one: say the Japanese for to remove, then the Japanese for to tie.
\end{quote}

\subsection*{You'll want to know: \ja{ゆびわ}}
\textbf{\ja{ゆびわ}} — \emph{yubiwa} — \textquotedblleft{}a ring (for a finger)\textquotedblright{}.

\begin{quote}
\textbf{\ja{ははの} \ja{ゆびわは} \ja{うつくしいです。}} — \emph{haha no yubiwa wa utsukushii desu} — \textquotedblleft{}My mother's ring is beautiful.\textquotedblright{}
\end{quote}

\begin{cousinweb}
\textbf{\ja{ゆびわ}} (\emph{yubiwa}) is native Japanese, built from \emph{yubi}, \textquotedblleft{}finger\textquotedblright{}, and \emph{wa}, \textquotedblleft{}ring, circle\textquotedblright{}, both inherited rather than borrowed (JMdict, KANJIDIC2).
\end{cousinweb}

\begin{grammarlens}[title={What you've built}]
One more word for everyday life at A2.
\end{grammarlens}

\subsection*{Your turn}
\begin{itemize}
\raggedright
  \item \emph{Say it:} \emph{yubiwa}
  \item \emph{Say it:} \emph{yubiwa}, once more
  \item \emph{Say it:} \emph{haha no yubiwa wa utsukushii desu} — My mother's ring is beautiful.
  \item \emph{Recall:} say the Japanese for to remove, then the Japanese for to tie, then say \emph{yubiwa} again
\end{itemize}

\begin{tcolorbox}[breakable,colback=teal!4,colframe=teal!35!black,title={Before you move on}]
What does \ja{ゆびわ} mean? (\textquotedblleft{}A ring (for a finger)\textquotedblright{}.) Say it once more.
\end{tcolorbox}
